%%%%%%%%%%%%%%%%%%%%%%%%%%%%%%%%%%%%%%%%%%%%%%%%%
% DOCUMENT SETUP: Do not make any changes to this section.
%%%%%%%%%%%%%%%%%%%%%%%%%%%%%%%%%%%%%%%%%%%%%%%%%

\documentclass[a4paper,10.5pt,twoside]{article}
\hyphenpenalty=8000
\textwidth=125mm
\textheight=200mm

\usepackage{fancyhdr}
\pagestyle{fancy}
\fancyhead{}
\fancyfoot{}
\raggedbottom

\usepackage{xurl}
\usepackage{graphicx}
\usepackage{alltt}
\usepackage{subcaption}
\usepackage{amsmath}
\usepackage[hidelinks]{hyperref}
\urlstyle{same}
\usepackage[T1]{fontenc}
\usepackage[utf8]{inputenc}
\usepackage{lmodern}
\usepackage{csquotes}
\usepackage{comment}
\usepackage{float}
\usepackage{xcolor}
\usepackage[a4paper, top=3cm, bottom=2cm, left=3cm, right=3cm]{geometry}
\usepackage[english]{babel}

\pagenumbering{arabic}
\setcounter{page}{1}

\hypersetup{
    colorlinks=true,
    linkcolor=blue,
    citecolor=blue,
    urlcolor=blue
}

%%%%%%%%%%%%%%%%%%%%%%%%%%%%%%%%%%%%%%%%%%%%%%%%%
\begin{document}

\fancyhead[LE]{\thepage\ \ \ MBD Midterm Report}
\fancyhead[RO]{Multimodal Misinformation Detection\ \ \ \thepage}

\begin{center}
{\Large \textbf{Midterm Report: Multimodal Image--Text Misinformation Detection}}\\[12pt]
{\normalsize
Ishimwe Karekezi Guy Gael\\
Lynne Chepkwony\\
Emile Lucky Muhigira
}
\end{center}

\section{Executive Summary}
We built a working multimodal misinformation detector using CLIP embeddings \cite{radford2021clip} and a Logistic Regression classifier.
On the held-out in-domain test split, we obtained \textbf{85\% accuracy} and \textbf{ROC-AUC = 0.9237}.
The main limitation is training size: while the raw source table has more than 500,000 rows, the final trainable subset is 9,979 samples after sampling, image availability checks, and preprocessing.

\section{Data Pipeline and Coverage}
\subsection{Dataset Description (Fakeddit)}
Fakeddit is a large-scale multimodal fake-news dataset built from Reddit posts, where each sample includes text, image metadata, and misinformation-oriented labels.
It provides both fine-grained and coarse-grained label settings; in this project we use the binary (\texttt{2-way}) setting for misinformation risk classification.
Fakeddit is suitable for this task because it contains naturally occurring image-text pairings and supports multimodal consistency analysis at scale.
In this report, label mapping is defined as: \textbf{0 = likely consistent (non-misinformation)} and \textbf{1 = likely misinformation}.

\subsection{Pipeline and Coverage}
Our workflow started from the Fakeddit multimodal table \cite{nakamura2019fakeddit} and then filtered to usable image-text pairs.

\begin{figure}[H]
\centering
\includegraphics[width=0.95\linewidth]{figures/data_pipeline_counts.png}
\caption{Dataset reduction through the preprocessing pipeline.}
\label{fig:data_pipeline}
\end{figure}

Final coverage details are as follows:
\begin{itemize}
    \item Raw multimodal train table: 564,000 rows
    \item Non-null image+text rows: 562,466
    \item Working sample: 15,000
    \item Rows with recovered images: 10,025 (66.83\% of sampled set)
    \item Final CLIP-ready set: 9,979
\end{itemize}

\begin{figure}[H]
\centering
\includegraphics[width=0.7\linewidth]{figures/class_distribution.png}
\caption{Class distribution in the final modeling table (\(n=9,979\)).}
\label{fig:class_dist}
\end{figure}

\section{Exploratory Data Analysis}
\subsection{Sample of the final modeling dataset}
\begin{table}[H]
\centering
\caption{Example rows after preprocessing.}
\label{tab:sample_rows}
\begin{tabular}{|l|p{6.6cm}|c|c|}
\hline
\textbf{id} & \textbf{text} & \textbf{label} & \textbf{text\_len} \\
\hline
2cdqb2 & a terrifying and chocolatey cookie monster & 0 & 42 \\
\hline
axhmr4 & the synagogue cared for by muslims bbc reel & 1 & 43 \\
\hline
b25ywe & the flight of this plane & 1 & 24 \\
\hline
2umqj8 & mr butterface and his little brother & 0 & 36 \\
\hline
1kog6d & this pylon looks like a sleeping owl & 0 & 36 \\
\hline
\end{tabular}
\end{table}

\subsection{Data quality checks}
The following EDA checks were completed:
\begin{itemize}
    \item Missing values for \texttt{id}, \texttt{text}, \texttt{label}, \texttt{text\_len}: 0 each.
    \item Duplicate text rows: 31.
    \item Duplicate IDs: 0.
    \item Duplicate (\texttt{id}, \texttt{text}) pairs: 0.
\end{itemize}

\subsection{Text-length behavior}
\begin{table}[H]
\centering
\caption{Text length summary from EDA (\(n=9,979\)).}
\label{tab:textlen_summary}
\begin{tabular}{|l|r|}
\hline
\textbf{Statistic} & \textbf{Value} \\
\hline
Count & 9,979 \\
Mean & 50.94 \\
Std & 32.40 \\
Min & 5 \\
25\% & 28 \\
50\% & 45 \\
75\% & 66 \\
Max & 294 \\
\hline
\end{tabular}
\end{table}

\subsection{Top terms by class}
Top tokens extracted during EDA:
\begin{itemize}
    \item Label 0: \texttt{the, this, and, for, with, from, you, his, that, new, like, happy, looks}.
    \item Label 1: \texttt{the, this, and, for, with, from, has, that, like, after, man, dog, found}.
\end{itemize}

\section{Full Statistical Findings}
\subsection{Baseline Logistic Regression}
Held-out metrics (test \(n=2495\)):
\begin{itemize}
    \item Accuracy: 0.85
    \item ROC-AUC: 0.9237
    \item Class 0: precision 0.82, recall 0.84, F1 0.83
    \item Class 1: precision 0.87, recall 0.85, F1 0.86
\end{itemize}

\subsection{Model Comparison}
\begin{figure}[H]
\centering
\includegraphics[width=0.95\linewidth]{figures/model_comparison_rocauc.png}
\caption{Held-out ROC-AUC comparison across candidate models.}
\label{fig:model_compare}
\end{figure}

The logistic baseline achieved the strongest ROC-AUC among tested models.

\subsection{Stability Check}
\begin{figure}[H]
\centering
\includegraphics[width=0.9\linewidth]{figures/cv_stability.png}
\caption{5-fold CV stability for Logistic Regression (ROC-AUC and F1).}
\label{fig:cv_stability}
\end{figure}

Cross-validation scores were stable:
\begin{itemize}
    \item Mean ROC-AUC: \(0.9301 \pm 0.0047\)
    \item Mean F1: \(0.8689 \pm 0.0074\)
\end{itemize}

\section{Preliminary Discussion}
\subsection{What the results mean}
The approach based on semantic consistency signals from pretrained multimodal embeddings is effective for our current benchmark slice.
Also, a linear classifier on top of CLIP-derived features is competitive and efficient.

\subsection{What was surprising}
Tree and boosting alternatives underperformed the logistic baseline in ROC-AUC despite being more complex.
Calibration slightly improved Brier score but reduced ROC-AUC, so we retained the base logistic model.

\subsection{Interpretation of the 9k subset}
It is reasonable to expect potential improvement with more compute and larger subsets, but this is not guaranteed.
Performance at scale will still depend on label quality, data coverage, and domain mismatch.
Given enough compute, training on a larger portion of the available dataset is a key next step and is expected to improve robustness beyond the current 9,979-sample setting.
\subsection{African-context scope for this stage}
At this midterm stage, we report the current in-domain benchmark and EDA results.
Dedicated African-context collection/evaluation is planned for a later semester milestone.

\section{Draft Related Work}
Prior work has shown multimodal gains for fake news detection, including SpotFake \cite{khattar2020spotfake}.
Our method follows lightweight transfer using CLIP \cite{radford2021clip} instead of end-to-end heavy architectures.
For data ecosystem grounding, we position our work alongside Fakeddit \cite{nakamura2019fakeddit}, MM-COVID \cite{li2020mmcovid}, and VisualNews \cite{liu2021visualnews}.
Out-of-context and transfer challenges are also highlighted in NewsCLIPpings \cite{luo2021newsclippings}.

\section{Limitations}
\begin{itemize}
    \item Main limitation: training subset size (9,979) relative to raw data availability.
    \item Data loss due to missing/unavailable images during local recovery.
    \item Current results are in-domain only; broader context transfer is not evaluated in this report version.
    \item African-context evaluation is deferred to later milestones.
\end{itemize}

\section{Conclusion}
At midterm, we have a working prototype that can be live-demonstrated.
Results are promising in-domain (\(\approx 85\%\) accuracy), and the next milestone focus is larger-scale training plus African-context evaluation.

\begin{thebibliography}{9}

\bibitem{radford2021clip}
Radford, A., Kim, J. W., Hallacy, C., et al. (2021).
Learning Transferable Visual Models from Natural Language Supervision.
\textit{Proceedings of the 38th International Conference on Machine Learning}.

\bibitem{nakamura2019fakeddit}
Nakamura, K., Levy, S., and Wang, W. Y. (2019).
r/Fakeddit: A New Multimodal Benchmark Dataset for Fine-grained Fake News Detection.
\textit{arXiv preprint arXiv:1911.03854}.

\bibitem{khattar2020spotfake}
Khattar, D., Goud, J. S., Gupta, S., and Varma, V. (2020).
SpotFake: A Multimodal Framework for Fake News Detection.
\textit{Proceedings of the IEEE International Conference on Big Data}.

\bibitem{li2020mmcovid}
Li, Y., Jiang, B., Shu, K., and Liu, H. (2020).
MM-COVID: A Multilingual and Multimodal Data Repository for Combating COVID-19 Disinformation.
\textit{arXiv preprint arXiv:2011.04088}.

\bibitem{liu2021visualnews}
Liu, F., Shi, J., and Liu, Y. (2021).
VisualNews: A Large-Scale Dataset for Image--Caption Analysis in News.
\textit{arXiv preprint arXiv:2101.01790}.

\bibitem{luo2021newsclippings}
Luo, G., Darrell, T., and Rohrbach, A. (2021).
NewsCLIPpings: Automatic Generation of Out-of-Context Multimodal Media.
\textit{Proceedings of the 2021 Conference on Empirical Methods in Natural Language Processing}.

\end{thebibliography}

\end{document}
